\subsection{Example: Counting words}

\begin{frame}[t,fragile]{Example: Counting words}
\begin{itemize}
\item Counting words in a text:
    \begin{itemize}
        \item An empty text has 0 words.
        \item A word starts with a non-letter character followed by a letter character.
        \item If text starts with a letter character, it starts with a word.
    \end{itemize}
\end{itemize}

\begin{block}{Checking for a new word}
\begin{lstlisting}
  size_t is_letter(char c) {
    return static_cast<std::size_t>(
        std::isalpha(static_cast<int>(static_cast<unsigned char>(c))) != 0);
  }

  size_t is_new_word(char c1, char c2) {
    return static_cast<std::size_t>((is_letter(c1) == 0) and (is_letter(c2) != 0));
  }
\end{lstlisting}
\end{block}

\end{frame}

\begin{frame}[t,fragile]{Counting words: Naive approach}
    \begin{itemize}
        \item An explicit loop checking consecutive characters.
        \item The first character is a special case.
    \end{itemize}
\begin{block}{Explicit loop}
\begin{lstlisting}
  size_t count_classic(std::string_view text) {
    if (text.empty()) { return 0; }
    size_t result = 0;
    for (size_t i = 1; i < text.size(); ++i) {
      if (is_new_word(text[i - 1], text[i]) != 0) { ++result; }
    }
    return result + is_letter(text[0]);
  }
\end{lstlisting}
\end{block}
\end{frame}

\begin{frame}[t,fragile]{Using an algorithm: the map-reduce approach}
    \begin{itemize}
        \item Map: check for new words in pairs of characters.
          \begin{itemize}
              \item Sequences: $(t_0, \ldots, t_{n-1})$ and $(t_1, \ldots, t_n)$.
              \item Mapping: \cppcode{is_new_word(c1,c2)} applied to pairs of characters.
          \end{itemize}
        \item Reduce: \cppcode{std::plus} for adding results of the mapping.
    \end{itemize}

\begin{block}{Map-reduce approach}
\begin{lstlisting}
  size_t count_classic_stl(std::string_view text) {
    if (text.empty()) { return 0; }
    auto result =
        std::transform_reduce(std::begin(text), std::prev(std::end(text)), 
        std::next(std::begin(text)), 
        0ZU,
        std::plus{}, 
        is_new_word);
    return result + is_letter(text[0]);
  }
\end{lstlisting}
\end{block}
\end{frame}

\begin{frame}[t,fragile]{Using a parallel map-reduce with C++17}
    \begin{itemize}
        \item The same map-reduce approach, but with a parallel execution policy.
    \end{itemize}

\begin{block}{Parallel map-reduce approach}
\begin{lstlisting}
  size_t count_parstl(std::string_view text) {
    if (text.empty()) { return 0; }
    auto result = std::transform_reduce(std::execution::par, 
        std::begin(text), std::prev(std::end(text)),
        std::next(std::begin(text)), 
        0ZU, 
        std::plus{},
        is_new_word);
    return result + is_letter(text[0]);
  }
\end{lstlisting}
\end{block}
\end{frame}

\begin{frame}[t,fragile]{Using ranges and views}
\begin{itemize}
  \item Steps:
    \begin{itemize}
      \item Generate a view with the $n-1$ first characters.
      \item Generate a view with the $n-1$ last characters.
      \item Zip the two views to get a view of pairs of characters.
      \item Count the number of pairs that start a new word.
    \end{itemize}
\end{itemize}

\begin{block}{Using ranges and views}
\begin{lstlisting}
  size_t count_ranges(std::string_view text) {
    if (text.empty()) { return 0; }
    auto compute = std::ranges::views::zip(
      text | std::ranges::views::take(std::ranges::size(text) - 1),
      text | std::ranges::views::drop(1));
    auto result = std::ranges::count_if(compute, [](auto char_pair) {
      auto [c1, c2] = char_pair;
      return is_new_word(c1, c2);
    });
    return result + is_letter(text[0]);
  }
\end{lstlisting}
\end{block}
\end{frame}

\begin{frame}[t,fragile]{Using parallel ranges and views with C++26}
\begin{itemize}
  \item The same approach as before, but with a parallel execution policy.
  \item \textmark{Note:} The parallel execution policy is not yet supported by all compilers.
    \begin{itemize}
      \item Experimental preview in OneAPI DPL.
    \end{itemize}
\end{itemize}

\begin{block}{Using parallel ranges and views}
\begin{lstlisting}
  size_t count_pranges(std::string_view text) {
    namespace stdx = oneapi::dpl;

    if (text.empty()) { return 0; }
    auto compute =
        std::ranges::views::zip(text | std::ranges::views::take(std::ranges::size(text) - 1),
                                text | std::ranges::views::drop(1));
    auto result = stdx::ranges::count_if(stdx::execution::par_unseq, compute, [](auto char_pair) {
      auto [c1, c2] = char_pair;
      return is_new_word(c1, c2);
    });
    return result + is_letter(text[0]);
  }
\end{lstlisting}
\end{block}
\end{frame}